\section{Experiments}
\label{sec:experiments}

% =====================================================================
% NOTE TO AUTHORS (remove before submission). Fill every "XX" with measured
% numbers. Items marked >>NEEDS<< require a decision/value from the experiments.
% AAAI rule: table captions go BELOW the table (already done here).
% =====================================================================

We design experiments to answer three research questions:
\textbf{(RQ1)} Does the escalation attack evade the state-of-the-art defense, and
does T-GAT restore detection where mean-pooling fails, without sacrificing
effectiveness on continuous attacks?
\textbf{(RQ2)} Which components of the temporal edge encoder drive the gain?
\textbf{(RQ3)} Does T-GAT generalize across LLM backbones, graph topologies, and
system scales?

\subsection{Experimental Setup}
\label{sec:exp_setup}

\paragraph{Attack scenarios.}
We evaluate on a suite of agent-attack scenarios spanning multiple attack
families. Following the multi-agent construction of~\citet{wang2025gsafeguard}, we
retain attack cases that fail against a single agent, so that any success within
the MAS arises purely from inter-agent propagation. The suite includes tool
attacks built from InjecAgent~\citep{zhan2024injecagent}, as well as additional
prompt-injection and memory-poisoning benchmarks. Orthogonal to the attack family,
we instantiate two \emph{temporal} attack regimes: \emph{(i)} a \emph{continuous}
attack, where compromised agents act maliciously from the first round and stay
consistent, and \emph{(ii)} our \emph{escalation} attack
(Section~\ref{sec:method_attack}), where compromised agents begin benign and
intensify over rounds.
% >>NEEDS<< finalize the exact benchmarks per family and cite them, e.g.,
% prompt injection on CSQA/MMLU/GSM8K and memory poisoning on PoisonRAG. If only a
% subset is run, narrow this sentence accordingly to avoid over-claiming.

\paragraph{MAS configuration.}
Each MAS has $N{=}8$ agents communicating for $K{=}3$ rounds, with the number of
attackers ranging in $\{1,2,3,4\}$ and edge density in
$\{0.2,0.4,0.6,0.8,1.0\}$. We evaluate four topologies: chain, tree, star, and
random.
% >>NEEDS<< confirm N, K, #attackers, densities, and #test samples actually used.

\paragraph{Baselines.}
We compare against \textbf{No Defense} and the state-of-the-art
\textbf{G-Safeguard}~\citep{wang2025gsafeguard}, the latter in its two aggregation
variants---\emph{mean-pooling} and \emph{last-turn}---which our method replaces
with a temporal edge encoder. All detectors share the same GNN backbone,
embeddings, and training data for a fair comparison.

\paragraph{Metrics.}
We report \emph{attack success rate} (ASR\,$\downarrow$), the fraction of honest
agents induced to execute the attacker's target tool, and \emph{detection}
quality via precision, recall, and F1 (\,$\uparrow$) on the attacker class.
% >>NEEDS<< confirm metric definitions match your eval (evaluate_output.py for ASR;
% F1/P/R from the detector). Add task-accuracy/utility if you report it.

\paragraph{Implementation details.}
We use MiniLM~\citep{wang2020minilm} ($D{=}384$) to embed utterances and
a 2-layer edge-featured GAT with hidden size 1024 and 8 attention heads; the
temporal edge encoder uses [\,XX\,] attention heads and hidden size [\,XX\,]. We
train with Adam (learning rate $10^{-3}$, weight decay $2{\times}10^{-4}$) for
[\,XX\,] epochs using temporal augmentation. The detector is trained on MAS
generated with [\,XX\,] (\emph{e.g.}, GPT-4o-mini) and transferred to other
backbones without retraining.
% >>NEEDS<< temporal_heads, temporal_hidden, epochs, training backbone, hardware.

\subsection{Main Results (RQ1)}
\label{sec:exp_main}

\paragraph{The escalation attack evades aggregation-based detection.}
Table~\ref{tab:main} shows that under the \emph{continuous} attack G-Safeguard
detects attackers reliably (F1 $=$ XX) and suppresses ASR from XX to XX. Under our
\emph{escalation} attack, however, its detection collapses to F1 $=$ XX and ASR
rises to XX, confirming that mean-/last-pooling is blind to attacks whose
malicious signal is back-loaded in time.
% claim-evidence: numbers from Table 1.

\paragraph{T-GAT defends both regimes.}
T-GAT recovers detection under escalation to F1 $=$ XX and lowers ASR to XX, while
matching G-Safeguard on the continuous attack (F1 $=$ XX, ASR $=$ XX). A single
T-GAT model thus handles both attack profiles, supporting our claim that modeling
\emph{how} agents change---not only their average content---is the discriminative
signal.

\begin{table}[t]
\centering
\setlength{\tabcolsep}{4pt}
\begin{tabular}{lcccc}
\toprule
& \multicolumn{2}{c}{\textbf{Continuous}} & \multicolumn{2}{c}{\textbf{Escalation}} \\
\cmidrule(lr){2-3}\cmidrule(lr){4-5}
Method & ASR\,$\downarrow$ & F1\,$\uparrow$ & ASR\,$\downarrow$ & F1\,$\uparrow$ \\
\midrule
No Defense              & XX & --  & XX & --  \\
G-Safeguard (last)      & XX & XX & XX & XX \\
G-Safeguard (mean)      & XX & XX & XX & XX \\
\textbf{T-GAT (Ours)}   & \textbf{XX} & \textbf{XX} & \textbf{XX} & \textbf{XX} \\
\bottomrule
\end{tabular}
\caption{Main results under continuous and escalation attacks, averaged over
topologies and attacker counts. ASR is attack success rate (lower is better); F1
is attacker-detection F1 (higher is better). Best in \textbf{bold}.}
% >>NEEDS<< averaging protocol (over which topologies/backbones) and the numbers.
\label{tab:main}
\end{table}

\subsection{Ablation Study (RQ2)}
\label{sec:exp_ablation}

\paragraph{Each temporal component contributes.}
Table~\ref{tab:ablation} isolates the parts of the temporal edge encoder under the
escalation attack. Removing the cross-turn statistics (variance and step-wise
change) causes the largest drop (F1 $-$XX), confirming they carry the escalation
signal; removing temporal self-attention or the temporal convolution costs F1
$-$XX and $-$XX, respectively. Replacing the encoder with plain mean-pooling
recovers the baseline's vulnerability, and disabling the fusion gate reduces
robustness on the continuous attack, showing the gate is what lets one model serve
both regimes.
% >>NEEDS<< ablation numbers.

\begin{table}[t]
\centering
\setlength{\tabcolsep}{5pt}
\begin{tabular}{lcc}
\toprule
Variant & ASR\,$\downarrow$ & F1\,$\uparrow$ \\
\midrule
\textbf{T-GAT (full)}            & \textbf{XX} & \textbf{XX} \\
\;\; w/o cross-turn statistics   & XX & XX \\
\;\; w/o temporal self-attention & XX & XX \\
\;\; w/o temporal convolution    & XX & XX \\
\;\; w/o fusion gate             & XX & XX \\
\;\; mean-pooling (= G-Safeguard) & XX & XX \\
\bottomrule
\end{tabular}
\caption{Ablation of the temporal edge encoder under the escalation attack. Each
row removes or replaces one component of the full model.}
% >>NEEDS<< numbers and confirm the exact set of ablated variants you run.
\label{tab:ablation}
\end{table}

\subsection{Generalization (RQ3)}
\label{sec:exp_generalize}

\paragraph{Across LLM backbones and topologies.}
Although T-GAT is trained on MAS from a single backbone, Table~\ref{tab:generalize}
shows it transfers to unseen backbones and topologies under the escalation attack,
reducing ASR by XX--XX\,points over G-Safeguard across
% >>NEEDS<< list the backbones (e.g., GPT-4o, Claude-3.5-Haiku, LLaMA-3.1-70B,
% DeepSeek-V3) and confirm the topology set.
all settings.

\paragraph{Across system scale.}
Because T-GAT inherits the inductive property of GNNs, a detector trained on
8-agent MAS applies directly to larger systems. Figure~\ref{fig:scale} reports ASR
as the number of agents grows to [\,XX\,], where T-GAT maintains [\,XX\,].
% >>NEEDS<< scale range and numbers; or drop this paragraph if not evaluated.

\begin{table}[t]
\centering
\setlength{\tabcolsep}{4pt}
\begin{tabular}{lccc}
\toprule
& No Defense & G-Safeguard & \textbf{T-GAT} \\
Backbone & ASR\,$\downarrow$ & ASR\,$\downarrow$ & ASR\,$\downarrow$ \\
\midrule
GPT-4o-mini       & XX & XX & \textbf{XX} \\
GPT-4o            & XX & XX & \textbf{XX} \\
Claude-3.5-Haiku  & XX & XX & \textbf{XX} \\
LLaMA-3.1-70B     & XX & XX & \textbf{XX} \\
DeepSeek-V3       & XX & XX & \textbf{XX} \\
\bottomrule
\end{tabular}
\caption{Generalization across LLM backbones under the escalation attack (ASR,
averaged over topologies). The detector is trained on a single backbone and
transferred without retraining.}
% >>NEEDS<< confirm backbone list and numbers.
\label{tab:generalize}
\end{table}

\subsection{Analysis}
\label{sec:exp_analysis}

\paragraph{Per-round dynamics.}
Figure~\ref{fig:dynamics} traces detection and ASR across dialogue rounds. Under
the escalation attack, G-Safeguard's per-round detection stays near chance while
contamination spreads, whereas T-GAT flags attackers by round [\,XX\,] and arrests
the spread, illustrating that early temporal cues are detectable even before the
attack reaches full intensity.
% >>NEEDS<< the per-round figure (fig:dynamics) and the round index XX.

% TODO(figure): \ref{fig:scale} scalability plot, \ref{fig:dynamics} per-round
% detection/ASR. Add when results are ready.
